Should a web agent choose how it sees a page --- text, pixels, or both --- per task, rather than once? Mostly not: the mode a router would send the hard tasks to is already best for all of them. The labels a router needs are produced at the agent's own success rate, so a weaker agent supplies fewer of them --- exactly where routing would be worth most. We measure six observation modes across eight site--model pairs (cells) on VisualWebArena and WebArena --- four of the modes screenshot-free, one pair differing only in the image. The modes are complementary: each solves tasks the others miss, they fail in structurally different ways, and which is best reverses between task sets. An oracle picking a winning mode per task therefore looks valuable, but that gain is inflated by run-to-run noise: rerunning one mode on the same tasks changes 5--14\% of outcomes on our hosted-backbone cells, so there a second run of a mode in hand buys about as much as a new mode. What survives is a cost bound: sending only the tasks no mode solves to the cheapest mode cuts cost 9.5--30.6\% in 8 of 8 cells at unchanged success. We test five routing policies --- picking the mode, deciding when to spend on the strong mode, a zero-cost rule read off the task text, a confidence cascade, and pooled cost tiers --- and none robustly beats fixing one well-chosen mode, or mixing fixed modes at random to meet a cost budget; the single exception is fragile and sits in our sparsest cell. We read the limit as one of today's agents rather than of routing: label supply and routing opportunity rise together (correlation 0.95 across cells, an association we did not manipulate), so a stronger agent could overturn it. We report the rerun bands and the resolution threshold they imply, so the same correction applies to any multi-arm ceiling claim.
